\chapter{Digital Audio Fundamentals}

\chapterline

We have spent the book so far building general machinery: signals as vectors, frequency as
rotation, the Fourier transform as projection, sampling and quantization as the bridge between
the continuous world and the computer. Now we point all of it at one of the richest and most
familiar signals there is -- \textbf{sound}. Every tool from the first half of the book now becomes
audible. By the end you will know, concretely and without hand-waving, how the pressure wave at
your eardrum becomes a list of integers in a file, why those integers are sampled $44{,}100$
times a second and stored in $16$ bits, what ``loud'' really means, what clipping is and how to
avoid it, and how to record and play a sound on a Raspberry Pi.

This is also the launchpad for the next several chapters. Once we agree on what a digital audio
signal \emph{is}, we can extract features from it (Chapter 27), shape it with effects and
synthesize it from scratch (Chapter 28), and analyze the signals of real instruments
(Chapter 29). And although our through-line project is a \emph{gesture} controller, audio is
both its eventual \emph{output} (it will emit MIDI that drives a synthesizer) and a possible
\emph{input} (a contact microphone on the wrist hears tendon snaps and taps). So sound matters
to us in two directions, and it pays to understand it from the ground up.

\section{From a pressure wave to a number}
\label{ch26:pressure}

Sound is a mechanical wave: a region of slightly compressed air chases a region of slightly
rarefied air outward from a vibrating source. What reaches your ear -- or a microphone -- is
the \emph{air pressure at one point} rising and falling over time. Strip away everything else
and a sound is exactly the object from Chapter 1: a function $p(t)$ giving pressure (really, the
tiny deviation from atmospheric pressure) at each instant.

A microphone is a transducer: it converts that pressure variation into a proportional voltage
$v(t)$. A dynamic mic moves a coil in a magnetic field; a condenser mic varies the gap of a
charged capacitor; a MEMS mic (the kind inside your phone, and the kind you will solder to a
Pi) flexes a micromachined membrane over a tiny back-plate. The physics differs, but the job is
identical: pressure in, voltage out, ideally with the voltage a clean scaled copy of the
pressure. From the signal-processing point of view we do not care which mic; we care that we now
have a continuous-time electrical signal $v(t)$ that \emph{is} the sound.

\begin{intuitionbox}
Picture the whole chain as a relay of copies of the same wiggle. Vibrating object $\to$ air
pressure $\to$ membrane motion $\to$ voltage $\to$ (after the steps below) a list of numbers
$\to$ (on playback) voltage $\to$ speaker cone motion $\to$ air pressure $\to$ your eardrum.
Every arrow is supposed to preserve the \emph{shape} of the wiggle. Digital audio is the art of
turning the wiggle into integers and back without your ear noticing the round trip.
\end{intuitionbox}

\begin{figure}[H]
\centering
\begin{tikzpicture}[
  node distance=4.5mm,
  box/.style={draw, rounded corners=2pt, minimum height=8mm, align=center,
    font=\scriptsize, inner sep=3pt},
  >={Stealth[length=2mm]},
]
\node[box, fill=sigblue!12] (src) {vibrating\\source};
\node[box, right=of src, fill=sigblue!12] (air1) {air\\pressure};
\node[box, right=of air1, fill=sigblue!12] (mic) {microphone\\(voltage)};
\node[box, right=of mic, fill=quotientorange!18] (adc) {sample +\\quantize};
\node[box, right=of adc, fill=quotientorange!28] (pcm) {PCM\\numbers};
\node[box, below=8mm of pcm, fill=quotientorange!18] (dac) {DAC\\(voltage)};
\node[box, left=of dac, fill=sigblue!12] (spk) {speaker\\cone};
\node[box, left=of spk, fill=sigblue!12] (air2) {air\\pressure};
\node[box, left=of air2, fill=targetgreen!16] (ear) {eardrum};
\draw[->] (src) -- (air1);
\draw[->] (air1) -- (mic);
\draw[->] (mic) -- (adc);
\draw[->] (adc) -- (pcm);
\draw[->] (pcm) -- (dac);
\draw[->] (dac) -- (spk);
\draw[->] (spk) -- (air2);
\draw[->] (air2) -- (ear);
\end{tikzpicture}
\caption{The audio chain, top row capture and bottom row playback. Every arrow should preserve the
shape of the wiggle; the two orange stages (sample plus quantize on the way in, reconstruct on the
way out) are where the continuous world meets the integers of this chapter.}
\label{ch26:fig-chain}
\end{figure}

Figure~\ref{ch26:fig-chain} lays out that relay end to end. To get from the continuous voltage
$v(t)$ to numbers we apply the two operations developed in Chapters 18 and 19, now in their audio
incarnation:

\begin{enumerate}
\item \textbf{Sampling} (discretize \emph{time}): measure $v(t)$ at regular instants
  $t = nT_s$, producing a sequence $v[n] = v(nT_s)$. The spacing $T_s$ is the sampling period;
  its reciprocal $f_s = 1/T_s$ is the \emph{sample rate}.
\item \textbf{Quantization} (discretize \emph{value}): snap each measured voltage to the nearest
  level on a finite grid, so it fits in a fixed number of bits. The number of bits is the
  \emph{bit depth}.
\end{enumerate}

The result -- a stream of quantized samples taken at a fixed rate -- is called \textbf{pulse-code
modulation}, PCM. It is the bedrock representation underneath WAV files, audio CDs, and the
buffers your sound card hands you. Everything fancier (MP3, AAC, Opus) is a clever way of storing
\emph{fewer} bits that decode \emph{back} to PCM before they reach a speaker. We will build PCM
from the geometry and treat compression as a later refinement.

\begin{figure}[H]
\centering
\begin{tikzpicture}
\begin{axis}[
  width=13cm, height=5.2cm,
  xlabel={time $t$ (ms)}, ylabel={value},
  xmin=0, xmax=2, ymin=-1.5, ymax=1.5,
  axis lines=left, enlargelimits=false,
  ytick={-1,-0.5,0,0.5,1},
  legend style={at={(0.99,1.25)}, anchor=north east, draw=none, fill=none, font=\small},
  legend cell align=left,
]
% continuous signal
\addplot[sigblue, thick, domain=0:2, samples=200] {sin(deg(2*pi*1.5*x)) * 0.9};
\addlegendentry{continuous $v(t)$}
% quantization grid lines (drawn first so stems sit on top)
\foreach \q in {-1.25,-1,-0.75,-0.5,-0.25,0,0.25,0.5,0.75,1,1.25} {
  \addplot[gray!35, dashed, samples=2, domain=0:2] {\q};
}
% sample instants as stems, snapped to quantization grid (steps of 0.25)
\addplot[specred, ycomb, mark=*, mark size=1.3pt, samples=21, domain=0:2]
  {0.25*round(4*0.9*sin(deg(2*pi*1.5*x)))};
\addlegendentry{quantized samples $v[n]$}
\end{axis}
\end{tikzpicture}
\caption{The two discretizations of PCM. Time is sampled (the dots occur only at multiples of
$T_s$); value is quantized (the dots land only on the dashed level grid). The smooth blue curve
is the original continuous signal; the red stems are what the computer actually stores.}
\label{ch26:fig-pcm}
\end{figure}

\begin{happeningbox}
Notice that two completely different things are happening in Figure~\ref{ch26:fig-pcm}, and it is
worth keeping them mentally separate forever. \emph{Sampling} is about the horizontal axis: we
only know the signal at certain instants. \emph{Quantization} is about the vertical axis: even at
those instants we only record a rounded value. Sampling errors show up as aliasing (Chapter 18);
quantization errors show up as a faint noise floor (Chapter 19). Audio quality is the project of
making both errors small enough that the ear cannot hear them.
\end{happeningbox}

\section{What the ear demands: the hearing range}
\label{ch26:hearing}

Audio is the one branch of signal processing where the specification is written by biology. We do
not need to capture sound perfectly; we need to capture it well enough to fool a particular
organ. So before choosing numbers we ask: what can a human ear actually detect?

\paragraph{Frequency range.} A healthy young ear responds to vibrations from roughly $20$~Hz to
$20{,}000$~Hz ($20$~kHz). Below $20$~Hz you feel rumble more than hear pitch; above $20$~kHz
(ultrasound) the cochlea simply does not respond. With age the upper limit retreats -- many
adults top out near $15$--$16$~kHz. This three-decade span (a factor of $1000$ in frequency) is
the single most important fact in audio engineering, because it sets how fast we must sample.

\paragraph{Amplitude range.} The ear's dynamic range is staggering. The quietest audible sound
pressure (the \emph{threshold of hearing}) and the loudest tolerable sound (the
\emph{threshold of pain}) differ by a factor of about $10^6$ in pressure amplitude -- a power
ratio of about $10^{12}$. No one wants to write ``the pressure was $0.0000063$ pascals,'' so
loudness is measured logarithmically, in \textbf{decibels} (dB). We will make decibels precise in
Section~\ref{ch26:loudness}; for now just register that the ear hears \emph{ratios}, not absolute
amounts, which is exactly why a logarithmic scale fits it.

\begin{figure}[H]
\centering
\begin{tikzpicture}
\begin{semilogxaxis}[
  width=13cm, height=5.4cm,
  xlabel={frequency (Hz, log scale)}, ylabel={loudness (dB SPL)},
  xmin=20, xmax=20000, ymin=-10, ymax=130,
  axis lines=left, enlargelimits=false,
  xtick={20,100,1000,10000,20000},
  xticklabels={20,100,1k,10k,20k},
  legend style={at={(0.96,0.85)}, anchor=north east, draw=gray!30, fill=white, fill opacity=0.9, text opacity=1, font=\footnotesize, rounded corners=2pt},
]
% threshold of hearing (schematic, U-shaped, lowest near 3-4 kHz)
\addplot[sigblue, thick, samples=80, domain=20:20000]
  {18 + 35*( (ln(x)-ln(3500))/2.4 )^2};
\addlegendentry{threshold of hearing}
% threshold of pain (roughly flat near 120 dB)
\addplot[specred, thick, samples=2, domain=20:20000] {120};
\addlegendentry{threshold of pain}
% mark the most sensitive region
\addplot[gray, dashed, samples=2] coordinates {(3500,-10) (3500,120)};
\node[anchor=south, font=\scriptsize, gray] at (axis cs:3500,-8) {most sensitive $\approx$ 2--5 kHz};
\end{semilogxaxis}
\end{tikzpicture}
\caption{The audible field (schematic). Sound is heard only between the two curves. The ear is
most sensitive around $2$--$5$~kHz (where speech consonants and a crying baby live), and far less
sensitive at the extremes of the range. The vertical gap is the available dynamic range at each
frequency.}
\label{ch26:fig-audiogram}
\end{figure}

\begin{connectionbox}
That the ear hears on a logarithmic scale in \emph{both} axes -- octaves in frequency (next
section) and decibels in amplitude -- follows the Weber--Fechner pattern in psychophysics, where perceived change tracks the \emph{logarithm} of a stimulus ($\dd(\ln x) = \dd x/x$, Chapter~42).
The recurring lesson: when a quantity spans many orders of magnitude and we care about relative
change, work in $\log$. We did it for the Bode plots of Chapter 16; we do it again here.
\end{connectionbox}

\section{Choosing the sample rate}
\label{ch26:samplerate}

The Nyquist--Shannon theorem (Chapter 18) is blunt: to capture all frequencies up to $f_{\max}$
without aliasing, you must sample at more than twice $f_{\max}$. Set $f_{\max} = 20$~kHz, double
it, and you get $40$~kHz as the bare minimum sample rate for full-range audio. Real systems push
a little above the minimum to leave room for a non-ideal anti-aliasing filter (which cannot drop
from passband to stopband infinitely fast). That margin is exactly why the standard rates sit a
bit over $40$~kHz.

\begin{whybox}
Why $44{,}100$~Hz of all numbers? History, not acoustics. Early digital audio masters were stored
on video tape (the only affordable medium with enough bandwidth in the late 1970s), using a PCM
adaptor that packed audio samples into the lines of a video frame. The arithmetic of NTSC and PAL
video lines -- so many usable lines per field, so many fields per second, so many samples per
line -- happened to land on $44{,}100$ samples per second for both formats. Sony and Philips
adopted it for the Compact Disc, and the world inherited it. It comfortably exceeds the
$40$~kHz Nyquist floor, giving the anti-alias filter a $2$~kHz transition band above $20$~kHz.
\end{whybox}

\begin{historybox}
The deeper root is older still. In 1928 Harry Nyquist, at Bell Labs, studied how many distinct
pulses a telegraph channel of a given bandwidth could carry; in 1937 Alec Reeves invented
pulse-code modulation (PCM) to protect telephone calls from cumulative analog line noise. In
1949 Claude Shannon generalized the mathematics into the sampling theorem we rely on today:
a bandlimited continuous signal is \emph{completely and losslessly reconstructible} from samples
taken faster than twice its highest frequency using sinc-interpolation filters. When digital
audio emerged commercially in the 1970s and 1980s with the Sony/Philips Compact Disc (1982),
it sparked intense debate: analog purists insisted that slicing continuous music into discrete
staircases destroyed its acoustic subtleties. Rigorous signal theory proved otherwise -- with
proper dither, quantization adds low-level white noise with zero harmonic distortion.
An active modern frontier is the development of low-latency neural audio codecs that compress
multichannel sound into fractions of a kilobit per second using deep generative representations.
\end{historybox}

The common rates you will meet:

\begin{center}
\begin{tabular}{@{}llp{8.2cm}@{}}
\toprule
Rate & Nyquist limit & Where it is used \\
\midrule
$8{,}000$~Hz   & $4$~kHz   & Telephone speech; intelligible but dull (no high consonant detail).\\
$16{,}000$~Hz  & $8$~kHz   & ``Wideband'' speech, voice assistants, many speech models.\\
$22{,}050$~Hz  & $\approx 11$~kHz & Half of CD rate; older games, low-bandwidth music.\\
$44{,}100$~Hz  & $\approx 22$~kHz & Audio CD; the default for music distribution.\\
$48{,}000$~Hz  & $24$~kHz  & Pro audio, video, film, and most computer sound hardware.\\
$96{,}000$~Hz  & $48$~kHz  & High-resolution recording and studio production.\\
\bottomrule
\end{tabular}
\end{center}

\begin{trapbox}
``More is better'' is a trap. Sampling music at $192$~kHz does not make it sound better to a human
ear -- there is nothing audible above $20$~kHz to capture, so you are storing ultrasonic noise at
four times the file size. Higher rates are genuinely useful \emph{during production} (they ease filter
design and reduce artifacts when pitch-shifting or applying nonlinear effects), then you render
down. For a final delivery file, $44.1$ or $48$~kHz is transparent. The same restraint applies to
our IMU and EMG work later: sample fast enough to clear Nyquist for the band you care about, plus
a margin, and not absurdly faster.
\end{trapbox}

\subsection*{A quick reality check by hand}

\begin{example}[Will it alias?]
\label{ch26:ex-alias}
You record at $f_s = 8000$~Hz (telephone rate) and someone whistles a pure tone at $6000$~Hz with
no anti-aliasing filter. What happens? The Nyquist frequency is $f_s/2 = 4000$~Hz. A $6000$~Hz
tone exceeds it, so it folds: the apparent frequency is
\[
  f_{\text{alias}} = \abs{f_s - f_{\text{in}}} = \abs{8000 - 6000} = 2000~\text{Hz}.
\]
The whistle is recorded as a $2$~kHz tone -- a genuinely wrong pitch. This is why every real ADC
front end has an analog low-pass (anti-alias) filter before the sampler, cutting everything above
$f_s/2$. The filter cannot be perfect, which is the whole reason we sampled a bit above the
$40$~kHz floor in the first place.
\end{example}

\section{Bit depth: how finely we measure each sample}
\label{ch26:bitdepth}

Sampling fixes \emph{when} we measure; bit depth fixes \emph{how precisely}. With $b$ bits per
sample we have $2^b$ available levels to represent the voltage. Audio almost always uses
\emph{signed} integers so that $0$ sits in the middle and the wave swings positive and negative:
$8$-bit audio runs $-128 \dots 127$, $16$-bit runs $-32768 \dots 32767$, and so on.

The key consequence is the \textbf{signal-to-quantization-noise ratio}. Rounding each sample to
the nearest level introduces an error of at most half a level. Treating that error as roughly
uniform noise (Chapter 19), the resulting SNR for a full-scale signal works out to the famous
rule of thumb
\[
  \mathrm{SNR}_{\text{dB}} \approx 6.02\,b + 1.76.
\]
Every bit you add buys about $6$~dB of dynamic range (Figure~\ref{ch26:fig-snr}). This single line
explains every standard bit depth at a glance:

\begin{center}
\begin{tabular}{@{}lll p{6.6cm}@{}}
\toprule
Bits $b$ & Levels $2^b$ & SNR $\approx 6b{+}1.76$ dB & Character \\
\midrule
$8$  & $256$        & $\approx 50$~dB  & Audibly hissy; old games, telephony.\\
$12$ & $4096$       & $\approx 74$~dB  & Decent; many low-cost ADCs and DACs.\\
$16$ & $65{,}536$   & $\approx 98$~dB  & CD quality; transparent for playback.\\
$24$ & $16.7$M      & $\approx 146$~dB & Studio recording; huge headroom margin.\\
\bottomrule
\end{tabular}
\end{center}

\begin{figure}[H]
\centering
\begin{tikzpicture}
\begin{axis}[
  width=12cm, height=5cm,
  xlabel={bit depth $b$}, ylabel={SNR (dB)},
  xmin=0, xmax=26, ymin=0, ymax=160,
  axis lines=left, enlargelimits=false,
  xtick={4,8,12,16,20,24}, ytick={0,40,80,120,160},
  legend style={at={(0.05,0.92)}, anchor=north west, draw=gray!30, fill=white, fill opacity=0.9, text opacity=1, font=\footnotesize, rounded corners=2pt},
]
\addplot[sigblue, thick, samples=2, domain=1:25] {6.02*x + 1.76};
\addlegendentry{$6.02\,b + 1.76$}
\addplot[only marks, specred, mark=*, mark size=1.6pt] coordinates {
  (8,49.9) (12,74.0) (16,98.1) (24,146.3)};
\addplot[gray, dashed, samples=2, domain=0:26] {120};
\node[anchor=south, font=\scriptsize, gray!90!black] at (axis cs:14,122) {ear's $\approx 120$ dB dynamic range};
\end{axis}
\end{tikzpicture}
\caption{Quantization SNR is a straight line in bit depth, about $6$~dB per bit. The dots mark the
standard depths from the table; $16$~bits already clears the ear's quiet-room-to-loud-passage gap,
and $24$~bits sits far above the ear's full dynamic range (dashed).}
\label{ch26:fig-snr}
\end{figure}

\begin{whybox}
Why is $16$ bits ``enough'' for a finished recording when the ear's dynamic range is about
$120$~dB? Because you never need to span the entire range \emph{in one recording at one volume
setting}. $16$ bits give $\approx 98$~dB of usable range -- more than the difference between a
quiet room and a loud passage of music -- and dithering (adding a tiny, sub-LSB random noise floor before
quantizing) turns correlated quantization distortion into a flat, practically inaudible hiss.
Studios record at $24$ bits not because the ear hears $146$~dB, but because
extra headroom means you can set levels conservatively, never clip, and still have noise far below
audibility after mixing.
\end{whybox}

\begin{connectionbox}
The $6.02b + 1.76$ formula is not specific to audio; it is the uniform-quantization SNR derived in
Chapter~19, applied directly to sound pressure samples. This
is the payoff of building general machinery first: a single derivation about rounding error now
tells us why CDs are $16$-bit and why a cheap $8$-bit sound chip hisses.
\end{connectionbox}

\subsection*{Integer versus floating point}

Inside a digital audio \emph{workstation}, samples are usually carried as $32$-bit
floating-point numbers in the range $[-1.0, +1.0]$, not integers. Float has two virtues for
\emph{processing}: its dynamic range is enormous (so intermediate sums never overflow), and
intermediate \emph{values above $1.0$ are not permanently lost} -- a signal can momentarily exceed
full scale inside a processing chain and be brought back down before the final integer
conversion. Integers, by contrast, have no room above the top level; exceed it and the sample is
permanently clamped. That clamping is clipping, the subject of the next section. The conversion
between the two worlds is a simple scale: $\text{float } 1.0 \leftrightarrow \text{int } 32767$
for $16$-bit.

\begin{recoverybox}
\textbf{Pause and think}: If a 16-bit ADC gives about 98 dB of dynamic range, why does recording at 24 bits ($146$ dB) matter if human hearing tops out around $120$ dB? In real recording environments, no microphone preamp is perfectly noise-free, and no performer plays at full scale at all times. A 24-bit recording lets you leave 18 dB of safety headroom so unpredictable loud transients never clip, while your quietest signals still sit comfortably $30$ dB above the digital quantization noise floor.
\end{recoverybox}

\section{Clipping, full scale, and headroom}
\label{ch26:clipping}

Every digital format has a hard ceiling: the largest value a sample can hold. We call it
\textbf{full scale} (often written $0$~dBFS, ``decibels relative to full scale''). When the
incoming signal tries to swing past full scale, the sample cannot follow -- it gets pinned at the
maximum. The rounded tops and bottoms of the waveform are \textbf{clipping}.

\begin{figure}[H]
\centering
\begin{tikzpicture}
\begin{axis}[
  width=13cm, height=5cm,
  xlabel={time}, ylabel={value},
  xmin=0, xmax=2, ymin=-1.5, ymax=1.5,
  axis lines=left, enlargelimits=false,
  ytick={-1,0,1}, yticklabels={$-1$ (full scale),$0$,$+1$ (full scale)},
  legend style={at={(0.99,1.28)}, anchor=north east, draw=none, fill=none, font=\small},
  legend cell align=left,
]
% original loud sine (amplitude 1.4, exceeds full scale)
\addplot[gray!55, thick, dashed, domain=0:2, samples=200] {1.4*sin(deg(2*pi*1.5*x))};
\addlegendentry{intended signal (too hot)}
% clipped version, clamped to [-1,1]
\addplot[specred, very thick, domain=0:2, samples=200]
  {max(-1, min(1, 1.4*sin(deg(2*pi*1.5*x))))};
\addlegendentry{what is recorded (clipped)}
% full scale lines
\addplot[sigblue!70, thick, dashed, domain=0:2, samples=2] {1};
\addplot[sigblue!70, thick, dashed, domain=0:2, samples=2] {-1};
\end{axis}
\end{tikzpicture}
\caption{Clipping. The intended sine (gray, dashed) swings beyond full scale, so its peaks are
flattened against the $\pm 1$ ceiling (red). The flat tops are new sharp corners in the waveform,
and sharp corners mean new high-frequency content -- harmonics that were not in the original.}
\label{ch26:fig-clip}
\end{figure}

\begin{happeningbox}
What does clipping \emph{do} to the sound, geometrically? Flattening the peaks turns a smooth
sinusoid into something closer to a square wave. From the Fourier series (Chapter 6) you know a
square wave is a sine plus odd harmonics at $3f, 5f, 7f, \dots$ (Figure~\ref{ch26:fig-clipspectrum}).
So clipping \emph{manufactures harmonics that were never played}. Mild clipping adds a gritty edge (guitar distortion pedals do
this deliberately and gently); hard clipping turns clean audio into a harsh buzz. In addition, those
new harmonics can exceed Nyquist and fold back as inharmonic, metallic tones -- clipping and
aliasing conspire to make digital overload sound especially nasty.
\end{happeningbox}

\begin{figure}[H]
\centering
\begin{tikzpicture}
\begin{axis}[
  width=12cm, height=4.8cm,
  xlabel={harmonic of $f_0$}, ylabel={relative level},
  xmin=0, xmax=10, ymin=0, ymax=1.15,
  axis lines=left, enlargelimits=false,
  xtick={1,3,5,7,9}, ytick={0,0.5,1},
  ybar, bar width=8pt,
]
% odd-harmonic amplitudes of a square wave: 1, 1/3, 1/5, 1/7, 1/9
\addplot[draw=specred, fill=specred!45] coordinates {
  (1,1) (3,0.3333) (5,0.2) (7,0.1429) (9,0.1111)};
\end{axis}
\end{tikzpicture}
\caption{Spectrum of a hard-clipped sine, which approaches a square wave: the fundamental plus odd
harmonics at $3f_0, 5f_0, 7f_0, \dots$ with amplitudes $1/3, 1/5, 1/7, \dots$. None of these were
in the original tone; clipping built them, and the high ones can fold back below Nyquist as
aliasing.}
\label{ch26:fig-clipspectrum}
\end{figure}

\textbf{Headroom} is the safety gap you deliberately leave between your typical signal peaks and
full scale, measured in dB. If your loudest peaks hit $-6$~dBFS, you have $6$~dB of headroom: a
sudden transient twice as loud as expected still fits. Recording engineers chase a balance -- too
little headroom risks clipping on an unexpected peak; too much wastes dynamic range and pushes the
quiet parts toward the quantization floor. A common practical target is peaks around $-12$ to
$-6$~dBFS while tracking, leaving comfortable margin.

\begin{example}[How much headroom is that?]
\label{ch26:ex-headroom}
You set your recording level so that the loudest expected peak reaches half of full scale, i.e.
amplitude $0.5$ on a $[-1,1]$ scale. How many dB of headroom is that? Headroom is the ratio of
full scale to peak, in dB:
\[
  \text{headroom} = 20\log_{10}\!\left(\frac{1.0}{0.5}\right) = 20\log_{10}(2) \approx 6.0~\text{dB}.
\]
So ``half of full scale'' is exactly $6$~dB of headroom -- another appearance of the
``factor of two is $6$~dB'' rule that also governs bit depth. A transient up to twice your
expected peak will just touch full scale without clipping.
\end{example}

\begin{trapbox}
Clipping is irreversible. Once peaks are flattened, the lost information is \emph{gone}; no plugin
can truly un-clip, only disguise it. This is the opposite of being a little too quiet, which you
can simply amplify later (at the cost of also lifting the noise floor). When in doubt, record
a touch quieter. The asymmetry -- recoverable quietness versus unrecoverable clipping -- is why
seasoned engineers leave generous headroom.
\end{trapbox}

\section{Loudness, decibels, and dBFS}
\label{ch26:loudness}

We have used decibels three times already; now let us pin them down. A decibel is a way of
expressing a \emph{ratio} of two quantities on a logarithmic scale. For \emph{power} (or energy,
or intensity):
\[
  L_{\text{dB}} = 10\log_{10}\!\left(\frac{P}{P_{\text{ref}}}\right).
\]
For \emph{amplitude} (voltage, pressure, sample value) -- which is the square root of power -- the
factor doubles to compensate for the square:
\[
  L_{\text{dB}} = 20\log_{10}\!\left(\frac{A}{A_{\text{ref}}}\right).
\]
The reference $A_{\text{ref}}$ is what makes the number meaningful. Two references dominate audio:

\begin{itemize}
\item \textbf{dB SPL} (sound pressure level): $A_{\text{ref}} = 20\,\mu\text{Pa}$, the nominal
  threshold of hearing. This measures sound \emph{in the air}. $0$~dB~SPL is barely audible;
  conversation is $\approx 60$~dB; a rock concert $\approx 110$~dB; pain $\approx 120$--$130$~dB.
\item \textbf{dBFS} (decibels relative to full scale): $A_{\text{ref}}$ is the maximum sample
  value. By construction $0$~dBFS is full scale and every real digital signal sits \emph{below}
  it, so dBFS values are negative (e.g. $-12$~dBFS).
\end{itemize}

\begin{intuitionbox}
The two ``$10$'' and ``$20$'' formulas are not two different definitions -- they are the same
definition seen through power versus amplitude. Power goes as amplitude squared, and
$\log(A^2) = 2\log A$, so the factor $10$ on power becomes $20$ on amplitude. Memorize one anchor
and you can rebuild the rest: \emph{a factor of $2$ in amplitude is $\approx 6$~dB; a factor of
$10$ in amplitude is exactly $20$~dB; a factor of $10$ in power is exactly $10$~dB.}
\end{intuitionbox}

\begin{example}[RMS and a sine's level]
\label{ch26:ex-rms}
``How loud is a signal'' is usually answered with its \emph{RMS} (root-mean-square) value, not its
peak, because RMS tracks delivered power and correlates with perceived loudness. For a sine of
amplitude $A$,
\[
  A_{\text{rms}} = \sqrt{\frac{1}{T}\int_0^T \big(A\sin(2\pi f t)\big)^2\,\dd t}
                 = \frac{A}{\sqrt{2}} \approx 0.707\,A,
\]
because the average of $\sin^2$ over a cycle is $1/2$. So a full-scale sine ($A = 1$) has
$A_{\text{rms}} \approx 0.707$, i.e. its RMS level is $20\log_{10}(0.707) \approx -3$~dBFS even
though its \emph{peak} is $0$~dBFS. This $3$~dB gap between peak and RMS for a sine -- its
``crest factor'' -- is why a meter showing peaks and a meter showing loudness never agree, and why
you must say \emph{which} you mean.
\end{example}

\begin{applicationbox}
Modern streaming services (Spotify, YouTube, Apple Music) ``loudness-normalize'' every track to a
common target measured in \textbf{LUFS} (loudness units relative to full scale), a refinement of
dBFS that weights frequencies the way the ear does and integrates over time. This is why a quiet
jazz track and a loud pop track play back at similar perceived volume without you touching the
knob. It also quietly ended the ``loudness war'' -- the decades-long arms race of squashing
dynamics to sound louder than the next track -- because over-compressed masters now simply get
turned \emph{down} to hit the same LUFS target, gaining nothing.
\end{applicationbox}

\section{The audio spectrum and octaves}
\label{ch26:octaves}

Because the ear hears frequency logarithmically, the natural unit of pitch is not the hertz but
the \textbf{octave}: a doubling of frequency. Two notes an octave apart ($f$ and $2f$) sound
``like the same note, higher'' -- the most consonant interval there is, because every harmonic of
the upper note coincides with a harmonic of the lower. The entire audible range, $20$~Hz to
$20$~kHz, spans about $10$ octaves ($2^{10} = 1024 \approx 1000$), which is why we drew the
audiogram on a log axis.

Western music subdivides each octave into $12$ equal-ratio steps (the equal-tempered semitone),
so the frequency ratio of one semitone is
\[
  r = 2^{1/12} \approx 1.0595.
\]
Standard tuning pins one note: A4 (the A above middle C) $= 440$~Hz. Every other note follows by
multiplying or dividing by $r$:
\[
  f(\text{note}) = 440 \cdot 2^{\,n/12},
\]
where $n$ is the number of semitones above ($n>0$) or below ($n<0$) A4. This is the exact map we
will need to \emph{emit} when the wrist controller plays an instrument: a detected gesture chooses
a MIDI note number, and MIDI note $m$ corresponds to $f = 440 \cdot 2^{(m-69)/12}$~Hz.

\begin{figure}[H]
\centering
\begin{tikzpicture}
\begin{semilogxaxis}[
  width=13cm, height=4.6cm,
  xlabel={frequency (Hz, log scale -- equal spacing = equal octaves)},
  ylabel={}, ytick=\empty,
  xmin=27, xmax=4400, ymin=0, ymax=1,
  axis lines=left, enlargelimits=false,
  xtick={27.5,55,110,220,440,880,1760,3520},
  xticklabels={A0\\27.5,A1\\55,A2\\110,A3\\220,A4\\440,A5\\880,A6\\1760,A7\\3520},
  x tick label style={align=center, font=\scriptsize},
]
% mark each A with a stem; equal spacing on log axis demonstrates octaves
\foreach \f in {27.5,55,110,220,440,880,1760,3520} {
  \addplot[sigblue, thick] coordinates {(\f,0) (\f,0.8)};
}
\node[anchor=south, font=\scriptsize] at (axis cs:440,0.82) {concert A = 440 Hz};
\end{semilogxaxis}
\end{tikzpicture}
\caption{Successive A notes, each an octave (a doubling) above the last. On a logarithmic axis the
octaves are \emph{evenly spaced} -- the same visual distance for every doubling -- which is exactly
how the ear perceives them.}
\label{ch26:fig-octaves}
\end{figure}

\begin{example}[From note to frequency]
\label{ch26:ex-note}
What frequency is middle C (C4)? Middle C is $9$ semitones below A4, so $n = -9$:
\[
  f = 440 \cdot 2^{-9/12} = 440 \cdot 2^{-0.75} \approx 440 \cdot 0.5946 \approx 261.6~\text{Hz}.
\]
And its MIDI note number is $60$, since A4 is MIDI $69$ and $69 - 9 = 60$. You can sanity-check
with the MIDI formula: $440 \cdot 2^{(60-69)/12} = 440 \cdot 2^{-9/12} \approx 261.6$~Hz. The two
formulas agree because they are the same equation written with different ``zero'' notes.
\end{example}

\paragraph{Rough map of the audio band.} It helps to carry a mental partition of the spectrum:
\emph{sub-bass} ($20$--$60$~Hz, felt more than heard), \emph{bass} ($60$--$250$~Hz, fundamental of
most instruments), \emph{midrange} ($250$~Hz--$4$~kHz, where speech and melody live and the ear is
most sensitive), \emph{presence} ($4$--$8$~kHz, consonant clarity and ``edge''), and
\emph{brilliance} ($8$--$20$~kHz, air and cymbal shimmer). Equalizers, the tone controls of audio,
simply boost or cut these regions; we will build the filters that do it in Chapter 28.

\section{Build it: record and play on a Raspberry Pi}
\label{ch26:build}

Enough theory -- let us make a Pi capture real PCM and play it back, so the numbers above stop
being abstract. A Raspberry Pi has \emph{no} analog audio input of its own, so for recording you
add a small USB sound card or an I2S MEMS microphone; its onboard $3.5$~mm jack (or HDMI) can play
\emph{out}. The software stack is ALSA underneath, usually wrapped by friendlier tools.

\begin{buildbox}
\textbf{Hardware.} A \$5 USB audio dongle (record + play) is the painless choice and shows up as
an ALSA card. For a tidier, all-digital path use an I2S MEMS mic such as the INMP441 or SPH0645,
wired to the Pi's I2S pins (BCLK, LRCLK, DOUT, plus 3V3 and GND) and enabled via a device-tree
overlay in \texttt{/boot/config.txt}. I2S skips analog wiring entirely: the mic does its own ADC
and streams PCM straight to the Pi.

\textbf{Find your devices.}
\begin{lstlisting}[language=bash]
arecord -l      # list capture (input) devices -> note card,device numbers
aplay   -l      # list playback (output) devices
\end{lstlisting}

\textbf{Record and play raw PCM.}
\begin{lstlisting}[language=bash]
# Record 5 s, 44.1 kHz, 16-bit signed little-endian, mono, from card 1 dev 0
arecord -D plughw:1,0 -f S16_LE -r 44100 -c 1 -d 5 take1.wav
# Play it back through the default output
aplay take1.wav
\end{lstlisting}
The flags \emph{are} this chapter: \texttt{-r 44100} is the sample rate, \texttt{S16\_LE} is the
$16$-bit signed PCM bit depth, \texttt{-c 1} is one (mono) channel.
\end{buildbox}

Once the bytes are on disk you can treat them as the integer signal we have been discussing. In
Python, \texttt{numpy} and \texttt{soundfile} (or the stdlib \texttt{wave} module) read a WAV file
straight into an array of samples:

\begin{lstlisting}[language=Python]
import numpy as np, soundfile as sf

x, fs = sf.read("take1.wav")          # x in [-1, 1] floats, fs = 44100
print("samples:", x.shape[0], " rate:", fs, " duration:", x.shape[0]/fs, "s")

peak = np.max(np.abs(x))              # peak amplitude, 0..1
rms  = np.sqrt(np.mean(x**2))         # RMS amplitude
dbfs_peak = 20*np.log10(peak + 1e-12)
dbfs_rms  = 20*np.log10(rms  + 1e-12)
print(f"peak {dbfs_peak:6.1f} dBFS   rms {dbfs_rms:6.1f} dBFS   "
      f"headroom {-dbfs_peak:4.1f} dB")

clipped = np.mean(np.abs(x) > 0.999)  # fraction of samples pinned at full scale
print(f"clipped fraction: {clipped:.4%}")
\end{lstlisting}

Run that on a recording where you spoke too loudly and the ``clipped fraction'' jumps off zero --
you have measured the very effect of Figure~\ref{ch26:fig-clip}. Run it on a quiet recording and
the headroom number balloons, confirming Example~\ref{ch26:ex-headroom}'s arithmetic with your own
voice.

\begin{labbox}
Record the same phrase three times: once speaking very softly, once at normal volume, once shouting into the
mic. Print peak dBFS, RMS dBFS, and clipped fraction for each. Predict before you run: which run
clips? Which has the most headroom? Then plot the waveforms and confirm by eye that the shout's
peaks are flat-topped while the quiet take rides far below full scale. You have just measured
headroom, clipping, and dynamic range directly on your own voice.
\end{labbox}

\begin{buildbox}
\textbf{Generate a tone to test playback} (no mic needed). This writes a clean A4 sine at
$-6$~dBFS so you can verify your output chain and hear the difference between clean and clipped:
\begin{lstlisting}[language=Python]
import numpy as np, soundfile as sf
fs, dur, f0 = 44100, 2.0, 440.0
t = np.arange(int(fs*dur)) / fs
clean = 0.5 * np.sin(2*np.pi*f0*t)               # peak 0.5 -> -6 dBFS headroom
hot   = np.clip(1.4*np.sin(2*np.pi*f0*t), -1, 1) # deliberately clipped
sf.write("clean_a4.wav", clean, fs, subtype="PCM_16")
sf.write("clipped_a4.wav", hot,  fs, subtype="PCM_16")
\end{lstlisting}
Play both through \texttt{aplay}. The clipped one buzzes -- those are the odd harmonics created
by flat-topping the sine wave, now clearly audible.
\end{buildbox}

\section{Channels, interleaving, and file formats}
\label{ch26:formats}

Real audio is rarely a single stream. \textbf{Stereo} carries two channels (left, right);
surround formats carry more. In a PCM buffer the samples are usually \emph{interleaved}: L, R, L,
R, $\dots$ in time order, so frame $n$ holds all channels for instant $n$ back-to-back. A
two-channel, $16$-bit, $44.1$~kHz stream therefore consumes
\[
  44100 \times 2 \times 2 = 176{,}400 \text{ bytes per second} \approx 10.1 \text{ MB/min},
\]
which is exactly the uncompressed bitrate of an audio CD. That number is why lossy codecs exist.

\begin{focusbox}
\textbf{The format zoo, demystified.}
\begin{itemize}
\item \textbf{WAV / AIFF} -- raw PCM in a thin container. Lossless, large, universally readable.
  Our default for analysis because the samples are \emph{exactly} what we recorded.
\item \textbf{FLAC / ALAC} -- \emph{lossless} compression: smaller files, bit-for-bit identical
  PCM on decode. Good for archiving.
\item \textbf{MP3 / AAC / Opus} -- \emph{lossy} compression: discards detail the ear is unlikely
  to notice (Chapter 33's psychoacoustic masking idea, applied to sound). Small files, but the
  PCM you get back is \emph{not} the PCM you put in. Never analyze fine spectral detail from a
  heavily compressed file.
\end{itemize}
\end{focusbox}

\begin{example}[Sizing a recording]
\label{ch26:ex-size}
How big is a $3$-minute stereo song at CD quality, uncompressed? Bytes per second is
$44100 \times 2~\text{ch} \times 2~\text{bytes} = 176{,}400$. Over $180$~s that is
$176{,}400 \times 180 = 31{,}752{,}000$ bytes $\approx 30.3$~MB. The same song as a $256$~kbit/s
AAC file is $256{,}000 \times 180 / 8 \approx 5.76$~MB -- about a fifth the size for
near-transparent quality. The factor-of-five-ish saving is the everyday payoff of lossy coding.
\end{example}

\section{Real limits, and the bridge to the wrist controller}
\label{ch26:claims}

\begin{goalbox}
Where does digital audio meet our gesture controller? In two solid places, and one tempting trap.

\textbf{Output (solid).} The controller's job is to \emph{make music}: a recognized gesture maps
to a MIDI note, and the note maps to a frequency by the equal-tempered formula of
Section~\ref{ch26:octaves}, which a synth turns into PCM. Every concept here -- sample rate, bit
depth, headroom, dBFS -- governs that output path. When we build the synthesizer in Chapter 28 we
will be writing exactly the PCM streams this chapter defined.

\textbf{Input (real but limited).} A \emph{contact} microphone or piezo pressed against the wrist
can pick up the faint acoustic ``clicks'' of tendons and the thud of a finger tap -- a field
called \emph{bioacoustic} sensing (Chapter 37). It is genuinely useful for detecting \emph{events}
(a tap happened; which of a few coarse gestures) using the audio pipeline of this chapter.

\textbf{The trap.} What a wrist microphone \emph{cannot} do is recover continuous, individual
finger \emph{joint angles}. The acoustic signal is a smeared, overlapping mixture; it does not
encode per-finger pose. Continuous finger tracking needs sensors \emph{on the fingers} (rings,
gloves) or, for muscle-level inference, EMG/FMG on the forearm -- and even those give limited,
calibrated estimates, not perfect joint angles. Treat wrist audio as an \emph{event detector},
never as a finger-pose sensor.
\end{goalbox}

\begin{aestheticbox}
Notice how directly human hearing dictates almost every engineering number in digital
audio. A $20$~Hz to $20$~kHz bandwidth sets the sample rate just above $40$~kHz; a $120$~dB
dynamic range sets the bit depth so quantization noise stays below the threshold of hearing;
logarithmic loudness and pitch perception give us decibels and octaves. Every standard audio
specification traces straight back to the physiology of the ear. Hold onto that habit -- letting the \emph{receiver}
define the requirements -- and it will guide every sensor choice in the chapters ahead, audio or
motion alike.
\end{aestheticbox}

We now have a firm grip on what a digital audio signal \emph{is}: PCM samples at a chosen rate and
bit depth, bounded by full scale, measured in decibels, carrying frequencies the ear can place on
a log scale of octaves. The next chapter turns these samples into \emph{features} -- energy, zero
crossings, spectral centroid, MFCCs -- the compact descriptions that let a machine (and our
gesture controller) tell one sound from another.

\section{Exercises}
\label{ch26:exercises}

\begin{exercise}
A mono recording is made at $f_s = 48{,}000$~Hz, $24$-bit, lasting $90$ seconds. (a) How many
samples are stored? (b) How many bytes (assume $3$ bytes per $24$-bit sample)? (c) What is the
Nyquist frequency, and is that enough to capture the full audible range?
\end{exercise}

\begin{exercise}
Using $\mathrm{SNR}_{\text{dB}} \approx 6.02\,b + 1.76$, find the approximate SNR for $b = 10$,
$14$, and $20$ bits. How many extra bits would you need to gain $30$~dB of dynamic range over a
$16$-bit recording, and why is that number what it is?
\end{exercise}

\begin{exercise}
A tone of $19$~kHz is recorded at $f_s = 32{,}000$~Hz with the anti-aliasing filter accidentally
switched off. (a) What is the Nyquist frequency? (b) Does the tone alias, and if so, to what
apparent frequency? (c) Is the aliased tone audible?
\end{exercise}

\begin{exercise}
You record a signal whose loudest peak reaches amplitude $0.25$ on a $[-1,1]$ scale. (a) What is
the peak level in dBFS? (b) How much headroom (in dB) do you have? (c) If you later amplify the
file by $9$~dB, what is the new peak amplitude, and have you clipped?
\end{exercise}

\begin{exercise}
Show, using the $20\log_{10}$ definition, that doubling a signal's amplitude raises its level by
about $6$~dB and that multiplying its amplitude by $10$ raises its level by exactly $20$~dB. Then
explain in one sentence why the power-based formula uses $10\log_{10}$ instead.
\end{exercise}

\begin{exercise}
(a) Compute the frequency of the note E5 (E above the treble staff), which is $7$ semitones above
A4. (b) Compute the frequency of A1, three octaves below A4. (c) What MIDI note number does each
correspond to, given A4 is MIDI $69$?
\end{exercise}

\begin{exercise}
A square wave at fundamental $f_0 = 5$~kHz is recorded at $f_s = 44{,}100$~Hz. Its harmonics fall
at $3f_0, 5f_0, 7f_0, \dots$. (a) List the harmonic frequencies up to $30$~kHz. (b) Which of them
exceed the Nyquist frequency and will alias? (c) Compute the aliased frequency of the first one
that folds. (This is why clipping -- which makes square-like waves -- can sound metallic.)
\end{exercise}

\begin{exercise}
A full-scale sine has peak $0$~dBFS. (a) Using $A_{\text{rms}} = A/\sqrt{2}$, what is its RMS
level in dBFS? (b) A full-scale square wave has $A_{\text{rms}} = A$ (it is always at $\pm A$).
What is \emph{its} RMS level in dBFS? (c) Explain why, at the same peak level, a square wave
sounds louder than a sine.
\end{exercise}

\begin{exercise}
Estimate the uncompressed storage for one hour of $5.1$ surround sound ($6$ channels) at
$48$~kHz, $24$-bit. Express the answer in gigabytes. Then estimate the same hour as a
$1.5$~Mbit/s lossy file and give the compression ratio.
\end{exercise}

\begin{exercise}
(Build.) Record a $5$-second clip with \texttt{arecord} at $44.1$~kHz/$16$-bit, then run the
Python snippet from Section~\ref{ch26:build}. Report peak dBFS, RMS dBFS, and clipped fraction.
Re-record while speaking much louder and report the same three numbers. Which changed most, and
does the clipped fraction behave the way Section~\ref{ch26:clipping} predicts?
\end{exercise}

\begin{exercise}
(Open-ended.) A product page advertises a wristband that ``listens to your hand with a
tiny microphone and reconstructs each finger's position in real time.'' Using
Section~\ref{ch26:claims}, write three sentences separating the plausible claim from the
implausible one, and name the sensor type that would actually be required for continuous
finger-pose tracking.
\end{exercise}
